\begin{proposition}[A bound on the compact modular region]\hypertarget{tag-001O}{}\label{ari:C.8.14}
\StatementRecord{ari:C.8.14}{proved}{001O}{\SeriesRef{ari:eq:cc-EZ}, \SeriesRef{ari:eq:cc-Fseries}, \SeriesRef{ari:eq:cc-ode}, \SeriesRef{ari:eq:cc-recurrence}, \SeriesRef{ari:eq:cc-power}, \SeriesRef{ari:eq:cc-finite-table}, \SeriesRef{ari:eq:cc-norms}, \SeriesRef{ari:eq:cc-Ffactored}, \SeriesRef{ari:eq:cc-scalar}, \SeriesRef{ari:C.8.12}}

Let $\chi=(\tau-i)/(\tau+i)$. On
$0<|\chi|\le2-\sqrt3$, the invariant $F$ satisfies $|F|<1$.
At the two equianharmonic points on the boundary its continuous
value is $3/4$.
\end{proposition}


\begin{proof}
\begin{enumerate}[label=\arabic*.,leftmargin=*]
\item Set
\[
 \begin{gathered}
 r_0=2-\sqrt3,\qquad A=E_4(i),\qquad k=\pi^2A,\\
 v=\pi\sqrt A\,\chi,\qquad z=v^2=k\chi^2,\\
 z_0=kr_0^2,\qquad c_0=1-\frac1{\sqrt3}.
 \end{gathered}
\]
In this proof $z$ is a coordinate on the modular disc. Define analytic functions $f_2,f_4,f_6$
by
\[
 \begin{gathered}
 \sqrt A f_2(v)=\frac{E_2(\tau)}{(1-\chi)^2}
                         -\frac3{\pi(1-\chi)},\\
 A f_4(v)=\frac{E_4(\tau)}{(1-\chi)^4},\qquad
 A^{3/2} f_6(v)=\frac{E_6(\tau)}{(1-\chi)^6}.
 \end{gathered}
\]
Ramanujan's differential identities become the rational system
\begin{equation}\hypertarget{tag-002T}{}\label{ari:eq:cc-ode}
 \begin{gathered}
 f_2'=\frac{f_4-f_2^2}{3},\qquad
 f_4'=\frac43(f_6-f_2f_4),\qquad
 f_6'=2(f_4^2-f_2f_6),\\
 (f_2,f_4,f_6)(0)=(0,1,0).
 \end{gathered}
\end{equation}
Uniqueness makes $f_4$ even and $f_2,f_6$ odd. Write
\[
 B(z)=f_4(v),\quad J(z)=\frac{f_6(v)}{2v},\quad
 D(z)=\frac{B(z)^3-4zJ(z)^2}{J(z)^{3/2}},\qquad
 H(z)=\sum_{n\ge0}\frac{-[z^n]D(z)}{16n^2-1}z^n.
\]
The branch has $J(0)^{1/2}=1$. The estimates below establish
its regularity on a larger disc, so these definitions are
unambiguous on $|z|\le z_0$.

\item Write $h=(1-\chi)(\mu-(1-\tau)/2)$. The Möbius chain rule and
\eqref{ari:eq:cc-EZ} give
\[
 h''(\chi)=-\frac{4\mu''(\tau)}{(1-\chi)^3}
       =-\frac{i}{2\sqrt{3\pi}}\chi^{-3/2}D(k\chi^2).
\]
The root expansion at $i$ gives the branch
\begin{equation}\hypertarget{tag-002R}{}\label{ari:eq:cc-h}
 h=\frac{2i}{\sqrt{3\pi}}\sqrt \chi\,H(z).
\end{equation}
Indeed differentiating this series twice gives precisely the
coefficient denominator $16n^2-1$ above. The two possible affine
homogeneous terms are absent because the half-period difference
has a Puiseux expansion consisting only of half-integral powers.
Changing the sign of the square root changes the signs of $h,p,R$
and leaves $F$ unchanged.

\item Put $r=|\chi|$. The definitions of $a$ and $R$ give exactly
\[
 R=\frac{1-\chi}{2i}
       \left(h'+\frac{\overline \chi h+\overline h}{1-r^2}\right),
 \qquad
 \frac p{g_2}=-\frac{i\sqrt \chi}{\sqrt{3\pi}(1-\chi)}
                         \frac{\sqrt J}{B}.
\]
Consequently
\begin{equation}\hypertarget{tag-002L}{}\label{ari:eq:cc-Fseries}
 F=\frac{\sqrt J}{B}
 \left\{H+4zH'+\frac{2r}{1-r^2}(rH-\overline H)\right\}.
\end{equation}
All conjugates in this display are values at the same point;
$H$ has real coefficients, so $\overline{H(z)}=H(\overline z)$.

\item The apparent pole at the equianharmonic points must be removed
before estimating this expression. We have $B(-z_0)=0$.
At these points $R=0$, since the expansion at $\rho$ gives
$\mu'=-1/3=a$. Thus, if
\[
 S=\sqrt J H,\qquad T=\sqrt J(H+4zH'),
\]
then $T(-z_0)=c_0S(-z_0)$. Define the analytic quotients
\[
 C(z)=\frac{B(z)}{1+z/z_0},\qquad
 U(z)=\frac{T(z)-c_0S(z)}{(1+z/z_0)C(z)},\qquad
 V(z)=\frac{\sqrt{J(z)}}{C(z)}.
\]
The estimates below prove $C\ne0$ on the needed larger disc.
For $r<r_0$, equation \eqref{ari:eq:cc-Fseries} is now
\begin{equation}\hypertarget{tag-002K}{}\label{ari:eq:cc-Ffactored}
 F=U+V\frac{dH+q(H-\overline H)}{1+z/z_0},
 \qquad
 d=\frac{2(r_0-r)}{(1+r_0)(1+r)},\quad q=\frac{2r}{1-r^2}.
\end{equation}
Here $q$ is only a real abbreviation within this display and
its estimates.

\item For an analytic
power series $K=\sum K_nz^n$, put
$\|K\|_s=\sum|K_n|s^n$. Then, for $0\le s\le z_0$,
\begin{equation}\hypertarget{tag-002S}{}\label{ari:eq:cc-norms}
 \begin{aligned}
 \|U\|_s&\le\frac1{\sqrt3}+\frac{13}{200}s,&
 \|V\|_s&\le1+\frac7{50}s,\\
 \|H\|_s&\le1+\frac{113}{1000}s,&
 \|H'\|_s&\le\frac{13}{100}.
 \end{aligned}
\end{equation}
In
\eqref{ari:eq:cc-ode}, write $f_2=\sum a_nv^n$,
$f_4=\sum b_nv^n$, $f_6=\sum c_nv^n$. Starting with
$(a_0,b_0,c_0)=(0,1,0)$ (only in this recurrence the symbol
$c_0$ denotes the zero initial coefficient), compute
\begin{equation}\hypertarget{tag-002Y}{}\label{ari:eq:cc-recurrence}
 \begin{aligned}
 a_{n+1}&=\frac{b_n-\sum_{j=0}^na_ja_{n-j}}{3(n+1)},\\
 b_{n+1}&=\frac{4(c_n-\sum_{j=0}^na_jb_{n-j})}{3(n+1)},\\
 c_{n+1}&=\frac{2(\sum_{j=0}^nb_jb_{n-j}-
                          \sum_{j=0}^na_jc_{n-j})}{n+1}.
 \end{aligned}
\end{equation}
Thus $B_n=b_{2n}$ and $J_n=c_{2n+1}/2$. For powers of a series
$K$ with $K_0=1$, the equally elementary recurrence is
\begin{equation}\hypertarget{tag-002V}{}\label{ari:eq:cc-power}
 [z^n]K^\alpha=\frac1n\sum_{j=1}^n
                 ((\alpha+1)j-n)K_j[z^{n-j}]K^\alpha.
\end{equation}
Together with $H_n=-D_n/(16n^2-1)$, these determine every
coefficient used. For example,
\[
 \begin{array}{c|rrrr}
 n&0&1&2&3\\ \hline
 B_n&1&10/9&5/27&4/81\\
 J_n&1&14/27&28/135&14/405\\
 H_n&1&13/135&-1999/153090&29299/28146690.
 \end{array}
\]
Division by $1+z/z_0$ uses $C_0=1$ and
$C_n=B_n-C_{n-1}/z_0$; the same recurrence applies to
$(T-c_0S)/(1+z/z_0)$. Division by $C$ uses the usual
convolution recurrence with constant coefficient one.

\item Here are sufficient rational enclosures for the only nonrational
constants in this arithmetic:
\begin{equation}\hypertarget{tag-002N}{}\label{ari:eq:cc-constants}
 \begin{gathered}
 14<k<72/5,\qquad
 1.0315619043558156493568263343177356527388253980203
 <z_0\\
 <1.0315619043558156493568263343177356527388253980205.
 \end{gathered}
\end{equation}
They follow from
$\pi=16\arctan(1/5)-4\arctan(1/239)$, the first 100 terms
of the alternating arctangent series, the exponential series
through degree 160, and
\[
 A=1+240\sum_{n=1}^{49}\frac{n^3e^{-2\pi n}}{1-e^{-2\pi n}}
       +\varepsilon_A,\qquad
 0<\varepsilon_A<
 \frac{240(4e^{-2\pi})^{50}}
 {(1-4e^{-2\pi})(1-e^{-2\pi})}.
\]
The last bound uses $n^3\le4^n$. Bounds for $\sqrt3$ follow
by squaring rational endpoints. All ensuing coefficient
comparisons are rational arithmetic using these enclosures.

\item To bound infinite tails, first consider $|z|=12/5$.
Here $|\chi|<5/12$, $\Im\tau>7/17$ and
$|e^{2\pi i\tau}|<1/13$. Lambert series give
\begin{equation}\hypertarget{tag-002U}{}\label{ari:eq:cc-outer}
 |B|<400,\qquad |J|<3000\qquad (|z|=12/5).
\end{equation}
For clarity the respective upper bounds before rounding are
\[
 \left(\frac{12}{7}\right)^4
 \left(1+240u\frac{1+4u+u^2}{(1-u)^5}\right)<400,
 \quad u=1/13,
\]
\[
 \frac13\left(\frac{12}{7}\right)^6
 \left(1+504u\frac{1+26u+66u^2+26u^3+u^4}{(1-u)^7}\right)<3000.
\]
We have used $A>1$ and $2\sqrt{12/5}>3$.
For $R=3/2$, let $\xi=5/8$ and retain coefficients through
degree 59. The tails of $B,J,C$ are respectively at most
\begin{equation}\hypertarget{tag-0031}{}\label{ari:eq:cc-tail60}
 \frac{400\xi^{60}}{1-\xi},\qquad
 \frac{3000\xi^{60}}{1-\xi},\qquad
 \frac{400z_0}{12/5-z_0}\frac{\xi^{60}}{1-\xi}.
\end{equation}
For the last bound, $B(-z_0)=0$ implies
$|C_n|\le400z_0(12/5)^{-n}/(12/5-z_0)$.
The finite sums from \eqref{ari:eq:cc-recurrence}--\eqref{ari:eq:cc-power}
with these tails give
\[
 \begin{array}{c|cccc}
 &\|B\|_{3/2}&\|J\|_{3/2}&\|C-1\|_{3/2}
                   &\min_{|z|\le3/2}\Re J(z)\\ \hline
 \text{bound}&<3276/1000&<2411/1000&<335/1000&>2049/10000.
 \end{array}
\]
For the last column, the lower bound used is exactly
\[
 1-J_2R^2-\frac{J_1^2}{8J_2}
       -\sum_{n=3}^{59}|J_n|R^n-
                   \frac{3000\xi^{60}}{1-\xi}.
\]
This is obtained by minimizing the quadratic
$\Re(1+J_1z+J_2z^2)$ on $|z|=R$; the minimum occurs at
$\cos\arg z=-J_1/(4J_2R)$. In particular, on $|z|\le3/2$,
\[
 |B|<33/10,\quad |J|<5/2,\quad |C|>3/5,\quad
 \Re J>1/5,\quad
 |D|\le\frac{(33/10)^3+6(5/2)^2}{(1/5)^{3/2}}<1000.
\]
Cauchy's inequality and the denominators in $H_n$ now give
\[
 |H|\le1+1000\sum_{n\ge1}\frac1{16n^2-1}
       =1+1000\left(\frac12-\frac\pi8\right)<110
                      \quad (|z|\le3/2).
\]
On $|z|\le7/5$, the same argument gives
\[
 |H+4zH'|\le1+1000\sum_{n\ge1}
                 \frac{(14/15)^n}{4n-1}
       \le1+\frac{1000}{3}\log15<910.
\]
It follows on $|z|=7/5$ that $|U|<10000$ and $|V|<3$;
indeed $|(1+z/z_0)C|>(7/(5z_0)-1)3/5$, and
$|\sqrt J|<8/5$.

\item Finally retain coefficients through degree 99 and put
$\eta=5z_0/7$. The tails in the coefficient norms for
$U,V,H,H'$ at $z_0$ are bounded respectively by
\begin{equation}\hypertarget{tag-0030}{}\label{ari:eq:cc-tail100}
 \frac{10000\eta^{100}}{1-\eta},\quad
 \frac{3\eta^{100}}{1-\eta},\quad
 \frac{110\eta^{100}}{1-\eta},\quad
 \frac{110\eta^{100}}{z_0}
 \left(\frac{100}{1-\eta}+\frac{\eta}{(1-\eta)^2}\right).
\end{equation}
The finite coefficient sums plus precisely these tails give
\begin{equation}\hypertarget{tag-002P}{}\label{ari:eq:cc-finite-table}
 \begin{array}{c|cccc}
 & (\|U\|_{z_0}-1/\sqrt3)/z_0
 & (\|V\|_{z_0}-1)/z_0
 & (\|H\|_{z_0}-1)/z_0 & \|H'\|_{z_0}\\ \hline
 \text{upper bound}&635/10000&1368/10000&1115/10000&1293/10000.
 \end{array}
\end{equation}
The recurrences, cutoffs and remainders above specify each
finite comparison by rational arithmetic.
Equations \eqref{ari:eq:cc-finite-table} imply
\eqref{ari:eq:cc-norms}, with room in each inequality.

\item For $s=|z|=kr^2$,
\[
 |H-\overline H|\le2|\Im z|\,\|H'\|_s,\qquad
 \frac{|\Im z|}{|1+z/z_0|}\le s.
\]
The second inequality follows by writing $z/z_0=ue^{i\phi}$
and using $(\cos\phi+u)^2\ge0$. The useful cancellation is
\[
 |dH+q(H-\overline H)|^2
      =d^2|H|^2+4q(d+q)(\Im H)^2.
\]
Also
\[
 \frac d{|1+z/z_0|}\le
 e(r):=\frac{c_0}{(1+r)(1+r/r_0)},\qquad
 q\le\kappa r,\quad d+q\le c_0+\kappa r^2,
 \quad\kappa=1+\frac2{\sqrt3}.
\]
The last bounds use $(1-r_0^2)^{-1}=1/2+1/\sqrt3$.
It follows from \eqref{ari:eq:cc-Ffactored}--\eqref{ari:eq:cc-norms},
using $k<72/5$, that
\begin{equation}\hypertarget{tag-002Z}{}\label{ari:eq:cc-scalar}
 |F|\le U_0(r)+V_0(r)
 \sqrt{e(r)^2H_0(r)^2+
       4\kappa r(c_0+\kappa r^2)(72/5)^2r^4(13/100)^2},
\end{equation}
where
\[
 U_0=\frac1{\sqrt3}+\frac{13}{200}\frac{72}{5}r^2,
 \quad V_0=1+\frac7{50}\frac{72}{5}r^2,
 \quad H_0=1+\frac{113}{1000}\frac{72}{5}r^2.
\]
To verify that the right side is less than one, let
\[
 M(r)=(1+r)^2(1+r/r_0)^2,
\]
\begin{equation}\hypertarget{tag-002M}{}\label{ari:eq:cc-P}
 \begin{aligned}
 P(r)={}&(1-U_0)^2M(r)\\
 &-V_0^2\left\{c_0^2H_0^2+
 4\kappa r(c_0+\kappa r^2)(72/5)^2r^4(13/100)^2M(r)\right\}.
 \end{aligned}
\end{equation}
This polynomial has $P(0)=0$. Write
\[
 \frac{P(r_0x)}{r_0x}
     =\sum_{j=0}^{14}b_j\binom{14}{j}x^j(1-x)^{14-j}.
\]
All coefficients belong to $\Q(\sqrt3)$. Explicit lower bounds,
obtained by expanding the displayed factored polynomial, are
\[
 \begin{array}{c|rrrrrrrr}
 j&0&1&2&3&4&5&6&7\\ \hline
 100b_j>&169&175&181&187&192&197&202&205\\[2pt]
 j&8&9&10&11&12&13&14&\\ \hline
 100b_j>&207&207&204&197&184&160&118&.
 \end{array}
\]
Explicitly, if $P(r)/r=\sum q_nr^n$,
then
\[
 b_j=\sum_{n=0}^j q_nr_0^n
                  \frac{\binom jn}{\binom{14}n}.
\]
For example $b_0=4-4\sqrt3/3$ and
$b_1=25357/13125-453\sqrt3/4375$. The rational enclosure
\[
 \begin{gathered}
 1.73205080756887729352744634150587236694280525381038<\sqrt3,\\
 \sqrt3<1.73205080756887729352744634150587236694280525381039.
 \end{gathered}
\]
verifies every entry by integer multiplication; its endpoints
are verified by squaring. Every Bernstein basis function is
nonnegative on $[0,1]$, so $P(r)>0$ for $0<r\le r_0$.
Since $U_0<1$ there, squaring \eqref{ari:eq:cc-scalar} is legitimate
and proves $|F|<1$. Formula \eqref{ari:eq:cc-Ffactored} applies
at $r=r_0$ except at $z=-z_0$ by continuity; at that remaining
point \ref{ari:C.8.12} gives $F=3/4$. This proves the assertion.
\end{enumerate}
\end{proof}